\scanpage{25}
divise \deleted{\unsure{cet}} \inserted{pseudo-}anneau à puissances divisées par l'idéal engendré par les
\[
(x.y)^{(n)}-x^n*y^{(n)}
\]
où le signe $*$ désigne provisoirement la multiplication de $\Gamma_A^+(J)$, le $.$ celle de $J$. On constate que les opérations $\pi^n$ passent au quotient \inserted{$J'$} pour donner une structure \inserted{de $A$-algèbre} à puissances divisées sur $J'$ [construction indépendante du fait que $J$ est un idéal de $A$]. On considère maintenant l'extension de $A/J$ par \inserted{la $A$-algèbre} $J'$ déduite de $J\longrightarrow J'$, i.e. le quotient
\[
A\times J'/\alpha(J^{\unsure{\flat|b}}),\quad\text{où}\qquad
\alpha:J^{\unsure{\flat|b}}\longrightarrow A\times J'
\]
est donné par
\[
\alpha(x)=(x,-x')
\]
où $x\sim x'$ est l'application canonique $J\longrightarrow J'!$

\deleted{On a, si $\deleted{\unsure{(a,X)}}\ (a,\unsure{\zeta|\xi})\in A\times J'$}

On \deleted{$(a,\unsure{\zeta|\xi})(x,-x')=ax,-\unsure{\zeta|\xi} x$} \unsure{je gue.}

\uline{N.B.} Supposons que $A$ soit une algèbre sur un anneau $\Lambda$. Alors le composé $\Lambda\longrightarrow A-\Gamma(A,J)$ définit évidemment sur $\Gamma(A,J)$ une structure de $\Lambda$-algèbre.
\mspage{24}
